% from ega3-6-fr.tex
\oldpage[III]{138}

façon générale les foncteurs de la forme précédente
$\mathscr G\mapsto\mathscr H_\bullet(\mathscr L_\bullet\otimes_Y\mathscr G)$
(où l'on fait en outre le cas échéant des hypothèses de finitude appropriées
sur les $\mathscr L_i$ ou les $\mathscr H_i(\mathscr L_\bullet)$) : c'est ce
qui est l'objet du § 7, dont la lecture, pour l'essentiel, est indépendante du
§ 6. Les propriétés les plus importantes de ces foncteurs concernent les
propriétés d'exactitude d'une composante $\mathscr T_i$ de
$\mathscr T_\bullet$; on donnera divers critères permettant d'établir de
telles propriétés; comme application, on obtiendra des conditions permettant
d'affirmer (avec les notations de (6.1.1)) que le foncteur
$\mathscr G\mapsto\mathrm R^n f_*(\mathscr F\otimes_{\mathscr O_Y}\mathscr G)$
est exact (ce qu'on exprimera en disant que $\mathscr F$ est
cohomologiquement plat sur $Y$ en dimension $n$). Une autre propriété
importante pour les composantes $\mathscr T_i$ de
$\mathscr T_\bullet(\mathscr G)=
\mathscr H_\bullet(\mathscr L_\bullet\otimes_Y\mathscr G)$ est une propriété
de semi-continuité de la fonction
$y\mapsto\dim_{k(y)}(T_i(k(y)))$; lorsque $\mathscr T_i$ est exact, cette
propriété est remplacée par une propriété de continuité, la réciproque étant
d'ailleurs vraie d'après Grauert lorsque $Y$ est réduit (7.8.4).
\end{env}

\begin{env}[6.1.3]
Dans les §§ 6 et 7, nous avons systématiquement fait usage de
l'hypercohomologie, en prenant partout comme arguments des complexes de
faisceaux au lieu de faisceaux, bien que la nécessité de ce point de vue
n'apparaisse que dans des chapitres ultérieurs. Le formalisme cohomologique
développé à cette occasion deviendra d'ailleurs plus transparent dans le
chapitre de ce Traité qui sera consacré à la mise au point d'une algèbre des
foncteurs cohomologiques de faisceaux cohérents, incluant le formalisme de la
dualité. Mais cela demandera des développements qui sortent du cadre du
présent chapitre.
\end{env}

\begin{env}[6.1.4]
Pour abréger, étant donnés deux complexes $K^\bullet$, $K'^\bullet$ dans une
catégorie abélienne $C$, nous dirons qu'un morphisme de complexes
$f:K^\bullet\to K'^\bullet$ est un homotopisme s'il existe un morphisme
$g:K'^\bullet\to K^\bullet$ tel que les morphismes composés $f\circ g$ et
$g\circ f$ soient tous deux homotopes à l'identité (par abus de langage,
lorsqu'il existe un tel homotopisme, on dira aussi que $K^\bullet$ et
$K'^\bullet$ sont homotopes). Lorsqu'on peut définir l'hypercohomologie d'un
foncteur covariant additif $T$ de $C$ dans une catégorie abélienne $C'$, par
rapport à un complexe de $C$ (0, 11.4.3), il est immédiat qu'un homotopisme
$K^\bullet\to K'^\bullet$ de complexes de $C$ définit canoniquement un
isomorphisme
\[
  \mathrm R^\bullet T(K^\bullet)\xrightarrow{\sim}
  \mathrm R^\bullet T(K'^\bullet)
\]
pour l'hypercohomologie (loc. cit.).
\end{env}

\subsection{Hypercohomologie des complexes de modules sur un préschéma}
\label{subsection:III.6.2.fr}

\begin{env}[6.2.1]
Soient $X$ un préschéma,
$\mathscr K^\bullet=(\mathscr K^i)_{i\in\mathbf Z}$ un complexe de
$\mathscr O_X$-Modules dont l'opérateur de dérivation est de degré $+1$.
Rappelons que pour tout morphisme $f:X\to Y$ de préschémas, on a défini
(0, 12.4.1) les $\mathscr O_Y$-Modules d'hypercohomologie
$\mathscr H^n(f,\mathscr K^\bullet)$ (aussi notés
$\mathscr H_f^n(\mathscr K^\bullet)$ ou
$\mathrm R^n f_*(\mathscr K^\bullet)$) pour tout $n\in\mathbf Z$;
l'hypercohomologie $\mathscr H^\bullet(f,\mathscr K^\bullet)$ est
l'aboutissement des deux foncteurs spectraux
$'\mathscr E(f,\mathscr K^\bullet)$ et
$''\mathscr E(f,\mathscr K^\bullet)$, dont les termes $\mathscr E_2$ sont
donnés par
\begin{align}
  '\mathscr E_2^{pq}
    &=\mathscr H^p\bigl(\mathscr H^q(f,\mathscr K^\bullet)\bigr),
    \tag{6.2.1.1}\label{III.6.2.1.1.fr}\\
  ''\mathscr E_2^{pq}
    &=\mathscr H^p\bigl(f,\mathscr H^q(\mathscr K^\bullet)\bigr)
      =\mathrm R^p f_*\bigl(\mathscr H^q(\mathscr K^\bullet)\bigr).
    \tag{6.2.1.2}\label{III.6.2.1.2.fr}
\end{align}

